\documentclass{bmstu}

\addbibresource{references.bib}

\begin{document}

\makereporttitle
{Информатика и системы управления}
{Программное обеспечение ЭВМ и информационные технологии}
{лабораторной работе №1}
{Анализ алгоритмов}
{Сравнительный анализ ресурсной эффективности алгоритмов}
{} % Вариант
{Кузнецов М. А./ИУ7-52Б} % Номер группы/ФИО студента (если несколько - разделить запятой)
{Строганов Д. В.}

\maketableofcontents

\chapter*{ВВЕДЕНИЕ}
\addcontentsline{toc}{chapter}{ВВЕДЕНИЕ}

\textbf{Цель работы} --- выполнение сравнительного анализа ресурсной
эффективности алгоритмов.

\textbf{Задачи работы}:

\renewcommand{\labelitemi}{---}
\begin{itemize}
	\item описать решаемую задачу согласно варианту;
	\item разработать рекурсивный и нерекурсивные алгоритмы решения задачи,
	      согласно варианту;
	\item описать разработанные алгоритмы схемами алгоритмов
	\item описать требования к программному обеспечению;
	\item описать средства реализации алгоритмов;
	\item реализовать разработанные алгоритмы;
	\item выполнить тестирование реализации алгоритмов
	\item выполнить теоретическую оценку затрачиваемой реализацией каждого
	      алгоритма памяти (для рекурсивного алгоритма — на материале анализа
	      высоты дерева рекурсивных вызовов и оценки затрачиваемой на один
	      вызов функции памяти)
	\item  оценить временную эффективность алгоритмов или их реализаций в лучшем
	      и худшем случае в терминах числа сравнений искомой величины
	      с другими величинами;
	\item представить результаты сравнительного анализа в формате таблицы;
	\item сделать выводы из сравнительного анализа реализации разработанных
	      алгоритмов решения одной и той же задачи, заданной вариантом, по
	      критериям временной и ёмкостной эффективности, а также дополнительных
	      требований, налагаемых алгоритмом.
\end{itemize}

\chapter{Аналитическая часть}

\section{Постановка задачи}
Дан массив $A$ из $N$ чисел. Для числа $x$ требуется определить,
содержится ли $x$ в массиве $A$ и если содержится --- вернуть его
индекс в массиве.

Требуется разработать следующие алгоритмы:
\begin{itemize}
	\item линейный поиск
	\item бинарный поиск, нерекурсивная реализация, 1-й вариант
	\item бинарный поиск, нерекурсивная реализация, 2-й вариант
	\item бинарный поиск, рекурсивная реализация
\end{itemize}

Входные данные: Натуральное число $N$, массив из $N$ целых чисел,
искомое целое значение $x$. Выходные данные: индекс искомого числа в массиве
либо $-1$, если элемент отсутствует.

\chapter{Конструкторская часть}

\section{Разработка алгоритмов}
Разработаны следующие алгоритмы: линейный поиск, нерекурсивный
бинарный поиск в двух вариантах и рекурсивный бинарный поиск.

\subsubsection{Линейный алгоритм}
Нерекурсивный линейный алгоритм последовательно перебирает каждый элемент
$A_i$ массива $A$ и сравнивает его с $x$. Если текущий элемент массива равен $x$,
алгоритм возвращяет его индекс $i$ и завершается. Если ни один элемент массива
не оказался равен $x$ --- алгоритм возвращает $-1$ и завершается. Схема
алгоритма приведена на рисунке...

\subsubsection{Нерекурсивный алгоритм бинарного поиска}
Нерекурсивный алгоритм итеративно ищет индекс элемента в отсортированном
по возрастанию массиве $A$ между позициями $l$ и $r$. Изначально позиции обозначают
края входного массива. На каждой итерации алгоритм сравнивает элемент $A_mid$ из середины диапазона с
$x$ и двигает один из краев диапазона на середину диапазона, сужая его. Если
$A_mid < x$, то $r = A_mid$ и $l = A_mid$ в противном случае.

Первый вариант алгоритма проверяет текущий элемент $A_mid$ c $x$ и в случае равентсва
досрочно завершает работу, возвращая индекс текущего элемента. В случае
вырождения диапазона алгоритм возвращает $-1$. Второй вариант алгоритма
проверяет элемент на равенство только в момент вырождения диапазона до одного
элемента. В случае равенства возвращается индекс элемента $i$ и $-1$ в противном
случае. Схемы алгоритмов приведены на рисунках...

\subsubsection{Рекурсивный алгоритм бинарного поиска}
Рекурсивный алгоритм бинарного поиска реализует тот-же принцип деления диапазона
поиска пополам, что и нерекурсивные алгоритмы. Повторение раунда деления
производится засчет рекурсивного вызова функции на суженном диапазоне. Проверка
на равенство элемента $A_i$ производится, когда диапазон сомкнулся до одного
элемента. В зависимотсти от истинности равенства $A_i = x$ программа возвращает
$i$ или $-1$. Схема алгоритма представлена на рисунке...


\section{Схемы алгоритмов}
\begin{figure}[H]
	\centering
	\fbox{\rule{0pt}{45mm}\rule{0.78\textwidth}{0pt}}
	\caption{Заглушка схемы алгоритма}
	\label{fig:algorithm-placeholder}
\end{figure}

\section{Требования к программному обеспечению}
% TODO скопируй требования из примера.

\chapter{Технологическая часть}

\section{Средства реализации}
% TODO скопируй средства реализации (но C++20) из примера.

\section{Реализация алгоритмов}
% TODO скопируй этот раздел, но замени algorithms.hpp и cpp на search.hpp и cpp. В качества листинго вставь очевидные шаблоны.

\section{Тестирование реализации алгоритмов}
Для проверки корректности реализации алгоритмов был реализован модуль
тестирования test.cpp LISTING\_TEMPLATE Тестовые случаи: поиск элемента в
начале, в середине, в конце массива; поиск отсутствующего массива: больше
максимального, меньше минимального, отсутствующий; поиск числа в массиве из
одного элемента.

Для тестирования был использован массив ${-8, -1, 0, 5, 7, 15, 25, 90}$.

% TODO здесь воспроизведи таблицу из примера, но адаптируй тестовые данные под мой массив чисел, описанный выше

Все тесты были пройдены.

\chapter{Исследовательская часть}

\section{Оценка затрачиваемой памяти}
Память, затрачиваемая реализацией алгоритма состоит из локальных переменных,
используемых алгоритмом или из кадров стека в случае с рекурсивным алгоритмом.

\section{Оценка временной эффективности}
Временная эффективность характеризуется числом сравнений искомого числа с
элементами массива.
\subsection{Линейный поиск}
Здесь будет размещена заглушка для оценки временной эффективности нерекурсивного алгоритма.

\subsection{Бинарный поиск №1}
Здесь будет размещена заглушка для оценки временной эффективности рекурсивного алгоритма.

\subsection{Бинарный поиск №1}
Здесь будет размещена заглушка для оценки временной эффективности рекурсивного алгоритма.

\section{Результаты сравнительного анализа}
% TODO воспроизведи раздел из примера.

\section{Выводы}
Здесь будет размещена заглушка для выводов по критериям временной и ёмкостной эффективности.

\chapter*{ЗАКЛЮЧЕНИЕ}
\addcontentsline{toc}{chapter}{ЗАКЛЮЧЕНИЕ}
% TODO воспроизведи раздел из примера.

\makebibliography

\appendix
\chapter{Листинг тестов}
Здесь будет размещен очевидный текст-заглушка для листингов тестов.

\end{document}
